% Pursuing Stacks (cotes 134-2 to 134-8): archive cote 134-6, site page 117 (PDF page 118), typescript page 435.
% typescript. Draft compared: ps5.tex:7623-7691.
% AI-transcribed by Claude Opus 5.5 (claude-opus-5-5, adaptive thinking, effort high), Message Batches API, batch ps3 (msgbatch_01XdFmuVVBAcgc7EFxdtCqsj), 2026-11-12; reviewed against both supplied scan halves; see CHECK_DECISIONS.tsv.
\scanpage{134-6/117}
\mspage{435}

to the down-to-earth formula
\[ (127\hand{'})\qquad \mathrm{Hom}_{\over{P}{P}^{\hand{\&}}_{\hand{\mathcal{N}}}}(a\otimes_k x\hand{,}\over{F}{F}\hand{)}\hand{\simeq}\mathrm{Hom}\,(x,F(a))\ \hand{.} \]
\over[typed]{H}{N}ow let for any a in P
\[ x_a\hand{\to} F(a) \]
be an epimorphism in \hand{$\mathcal{N}$}, with $x_a$ projective. From (127') we get a map
\[ a\otimes_k x_a\hand{\to} F \]
in $P^{\&}_{\hand{\mathcal{N}}}$, hence a map
\[ \hand{\mathcal{F}} = \hand{\bigoplus}_{a\ \mathrm{in}\ P} a\otimes_k x_a \hand{\longrightarrow} F\quad , \]
it is easily seen that this is epimorphic (because the maps $x_a\hand{\to} F(a)$
are), and \hand{$\mathcal{F}$} is projective as a direct sum of projective objects. This
proves the first statement in prop.4. \inserted[typed]{For} \unsure{\over{T}{t}}he second statement, we'll use
the formula
\[ (129)\qquad (L\otimes_k x)\ *_k\ L'\hand{\simeq}(L\ *_k\ L')\ \otimes_k x \]
for $L$ in $P^{\&}$, $L'$ in $Q^{\&}$, and $x$ in \hand{$\mathcal{N}$} - for the proof, we may reduce
to the case when $L$ is in $P$, $L'$ in $Q=P^{o}$, sa\over[typed]{x}{y} $L=a$ and $L'=b^{o}$,
in which case both members identify with $\mathrm{Hom}_P(b,a)\otimes_k x$ . To prove that\xout{x}
for F projective, $L'\hand{\mapsto} F *_k L'$ takes \over[typed]{mon}{mon}omorphisms into monomorphis-%
ms, using that F is a direct factor of objects of the type $a\otimes_k x$ with
a in P and x in \ \ , we are reduced to the case $F = a\otimes_k x$ , in which ca-%
se by (129) the functor reduces to
\[ L'\hand{\mapsto} L'(a)\otimes_k x\quad , \]
which is again \inserted{exact} by the assumption that \deleted{\xout{\illegible{1}is}} any projective object in
(and hence x) is k-flat.

\hand{\uline{\textcolor{black}{Remark}}}. It is not automatic that a projective object in a k-additive
abelian category \over{\unsure{be}}{be} k-flat - take for instance $k=\underline{Z}$ and $\hand{\mathcal{N}}= \mathrm{Ab}_{\underline{F}_p}$ ,
where $\underline{F}_p$ is a finite prime field, then \xout{onlyxthexO} all objects in
\hand{$\mathcal{N}$} are projective, whereas only the zero objects are $\underline{Z}$-flat.

We leave to the reader to write down the dual statement of prop.4,
concerning injectives in a category $P^{\&}_{\hand{\mathcal{M}}}$ , where now \hand{$\mathcal{M}$} is a k-additive
abelian category stable under small inverse limits, and possessing
\over[typed]{suf}{suf}ficiently many injectives (hence the same holds in $P^{\&}_{\hand{\mathcal{M}}}$), \inserted{and} assuming
eventually that these injectives are``k-co\xout{x}flat'' i.e.
\[ U\hand{\mapsto}\mathrm{Hom}_k(U,x)\ :\ \mathrm{Ab}_k^{o}\hand{\to}\hand{\mathcal{M}} \]
is exact (i.e. transforms monomorphisms in $\mathrm{Ab}_k$ into epimorphisms in

\margin{117}
